% !TeX root = tools-in-complex-dynamics.tex
\section{Tools in Complex Dynamics}


\subsection{Differential forms and currents}

The reference for this section is, Demally, holomorphic dynamics, Dinh-Sibony notes. \Yang{}
\cite{Dinh-Sibony-2005-introduction-to-currents}

Let \(M\) be a complex manifold of complex dimension \(d\).
Recall that we have the decomposition of the cotangent bundle:
\[ \Omega_{\sm}^{1,\bbC} \cong \Omega_{\sm}^{1,0} \oplus \Omega_{\sm}^{0,1} \cong \Omega_{\sm}^1 \oplus \overline{\Omega_{\sm}^1}. \]
Take exterior powers, we have \( \Omega_{\sm}^{k,\bbC} \cong \bigoplus_{p+q=k} \Omega_{\sm}^{p,q} \),
where \(\Omega_{\sm}^{p,q} = \bigwedge^p \Omega_{\sm}^{1,0} \otimes \bigwedge^q \Omega_{\sm}^{0,1}\).
We also use the notation
\[ \calA^{p,q} \coloneqq \Omega_{\sm}^{p,q}, \quad \calA^{k} \coloneqq \Omega_{\sm}^{k,\bbC}. \]

A reason to induce this strange sheaf \(\calA^{k} = \Omega_{\sm}^{k,\bbC}\) is to make sense of integration of top-degree forms.
For simplicity, assume that \(M\) is compact.
Let \(\omega \in \calA^{2d}(M)\) be a smooth complex-valued \(2d\)-form on \(M\).
Then its integration is well-defined in the smooth manifold sense:
\[ \int_M \omega \in \bbC. \]
However, in complex case, it is more natural to ``integral'' a holomorphic \(d\)-form on a \(d\)-dimensional complex manifold.
This does not make sense in the smooth manifold theory.
The solution is to associate a holomorphic \(d\)-form \(\eta \in \calA^{d,0}(M)\) with a smooth \(2d\)-form (\((d,d)\)-form) \(\omega = \eta \wedge \overline{\eta} \in \calA^{d,d}(M) \subset \calA^{2d}(M)\).

Another reason  is that \(\bigoplus_k \Omega_{\sm}^{k}\) is not closed under the exterior derivative \(\upd\), while \(\bigoplus_k \Omega_{\sm}^{k,\bbC}\) is.
Suppose that we have local holomorphic coordinates \((z_1, \ldots, z_d)\).
Recall that we have the exterior derivative
\[ \upd: \Omega_{\sm}^{1,0} \to \Omega_{\sm}^{2,\bbC}, \quad f d z_{i} \mapsto \sum_{j=1}^d \frac{\partial f}{\partial z_j} d z_j \wedge d z_{i} + \sum_{i=1}^d \frac{\partial f}{\partial \overline{z_j}} d \overline{z_j} \wedge d z_{i}. \]
On its conjugation, we have
\[ \upd: \Omega_{\sm}^{0,1} \to \Omega_{\sm}^{2,\bbC}, \quad f d \overline{z_{i}} \mapsto \sum_{j=1}^d \frac{\partial f}{\partial z_j} d z_j \wedge d \overline{z_{j}} + \sum_{j=1}^d \frac{\partial f}{\partial \overline{z_j}} d \overline{z_j} \wedge d \overline{z_{j}}. \]

Extending \(\upd\) by linearity and the Leibniz rule \(\upd(\alpha \wedge \beta) = \upd \alpha \wedge \beta + (-1)^{\deg \alpha} \alpha \wedge \upd \beta\), we get the exterior derivative
\[ \upd: \calA^k = \Omega_{\sm}^{k,\bbC} \to \calA^{k+1} = \Omega_{\sm}^{k+1,\bbC}, \]
which can be decomposed as \(\upd = \partial + \overline{\partial}\), where
\begin{align*}
	\partial: \calA^{p,q}            & \to \calA^{p+1,q}, \quad f d z_I \wedge d \overline{z_J} \mapsto \sum_{j=1}^d \frac{\partial f}{\partial z_j} d z_j \wedge d z_I \wedge d \overline{z_J},                       \\
	\overline{\partial}: \calA^{p,q} & \to \calA^{p,q+1}, \quad f d z_I \wedge d \overline{z_J} \mapsto \sum_{j=1}^d \frac{\partial f}{\partial \overline{z_j}} d \overline{z_j} \wedge d z_I \wedge d \overline{z_J}.
\end{align*}
In a diagram, we have:
\begin{center}
	\begin{tikzpicture}
		\matrix (m) [matrix of math nodes, column sep=2em, row sep=2em] {
			\bullet & \bullet       & \bullet       & \bullet   & \bullet     \\
			\bullet & \calA^{p,q+1} & \bullet       & \bullet   & \bullet     \\
			\bullet & \calA^{p,q}   & \calA^{p+1,q} & \bullet   & \bullet     \\
			\bullet & \bullet       & \bullet       & \bullet   & \bullet     \\
			        &               &               & \calA^{k} & \calA^{k+1} \\
		};
		% Draw a horizontal line before the row containing \Omega^k (above row 5)
		\draw (m-4-1.south) -- (m-4-5.south);
		% Draw a dashed line passing through \Omega^{p-1,q+1} and \Omega^k
		\draw[dashed] (m-2-1) -- (m-5-4);
		\draw[dashed] (m-1-1) -- (m-5-5);
		% Draw arrows for the operators
		\draw[->] (m-3-2) -- (m-3-3) node[midway, above] {\(\partial\)};
		\draw[->] (m-3-2) -- (m-2-2) node[midway, right] {\(\overline{\partial}\)};
		% \draw[->] (m-3-2) -- (m-4-4) node[midway, right] {\(\mu\)};
		% \draw[->] (m-3-2) -- (m-1-1) node[midway, above] {\(\overline{\mu}\)};
		\draw[->] (m-5-4) -- (m-5-5) node[midway, above] {\(\upd\)};
	\end{tikzpicture}
\end{center}

\begin{definition}\label{def:forms_with_compact_support}
	A differential form \(\omega \in \calA^{k}(M)\) is said to have \emph{compact support} if there exists a compact subset \(K \subset M\) such that \(\omega|_{M \setminus K} = 0\).
	The space of smooth complex-valued \(k\)-forms with compact support on \(M\) is denoted by \(\calA^{k}_c(M)\).
\end{definition}

\begin{definition}
	Given a smooth local coordinate chart \((U, (x_1, \ldots, x_{2d}))\), a compact subset \(K \subset U\), and a non-negative integer \(m\), writing a form $\omega = \sum_I u^I \upd x_I$ under the coordinate, we can define a seminorm
	\[ p_{U,K,m}(\omega) = \sup_{x \in K} \max_{I} \max_{|\alpha| \leq m} |D^{\alpha} u^I(x)|, \]
	where \(\alpha = (\alpha_1, \ldots, \alpha_{2d})\) is a multi-index, and \(D^{\alpha} = \frac{\partial^{|\alpha|}}{\partial x_1^{\alpha_1} \cdots \partial x_{2d}^{\alpha_{2d}}}\) in local coordinates, and
	where \(K\) runs over all compact subsets of \(M\) and \(m\) runs over all non-negative integers.
	We equip the spaces $\calA^k(M)$ (resp. $\calA^k_c(M)$) the weak topology induced by the seminorms $\{p_{U,K,m}\}_{U,K,m}$, and denote the resulted space by $\scrE^k(M)$ (resp. $\scrD^k(M)$).
	% $\scrE^{k}$ and $\scrD^{k}$ with the weak topology induced by these seminorms $\{p_{U,K,m}\}_{U,K,m}$.
\end{definition}

\begin{remark}
	Recall that the weak topology on \Yang{}
\end{remark}

Note that above definition works for any smooth manifold, not just complex manifolds.
The notions of currents was introduced by de Rham in the 1950s, and they provide a general framework for differential forms, cycles, measures, and distributions.

\begin{definition}
	The \emph{space of currents} of degree \(k\) or dimension $2d - k$ is defined as the topological dual of the space $\scrD^{2d-k}(M)$ of smooth \((2d-k)\)-forms with compact support:
	\[ \scrD'^k(M) \coloneqq \scrD'_{2d-k}(M) \coloneqq \scrD^{2d-k}(M)^\vee. \]
\end{definition}

\begin{remark}
	Recall that on a smooth manifold \(M\) of dimension \(n\), a \emph{distribution} is a continuous linear functional on the space of smooth functions with compact support, denoted by \(\scrD(M)\) or \(C_c^\infty(M)\).
	\Yang{}
\end{remark}

\begin{example}
	Any submanifold \(N \subset M\) of real dimension \(2d - k\) defines a current of degree \(k\) (or dimension $2d-k$) by integration:
	\[ T_N: \scrD^{2d-k}(M) \to \bbC, \quad \omega \mapsto \int_N \omega|_N. \]
	More generally, if \(N\) is a subvariety of complex dimension \(d - k\), then \(N\) defines a current of degree \(2k\) (or dimension $2d-2k$) by integration:
	\[ T_N: \scrD^{2d-2k}(M) \to \bbC, \quad \omega \mapsto \int_{N_{reg}} \omega|_{N_{reg}}, \]
	where \(N_{reg}\) is the regular part of \(N\).
	Since \(N\) is a complex subvariety, the singular part \(N_{\sing} = N \setminus N_{reg}\) is a proper subvariety of \(N\) of complex dimension at most \(d - k - 1\), hence the integral makes sense.
	Moreover, for any $k$-cycle \(Z = \sum_i n_i Z_i \in Z_k(M)\), where \(Z_i\) are irreducible subvarieties of complex dimension \(d - k\) and \(n_i \in \bbZ\), we can associate a current of degree \(2k\) to \(Z\) by linearity.
\end{example}

\begin{example}
	Any smooth \(k\)-form \(\omega \in \scrD^k(M)\) defines a current of degree \(k\) via the pairing
	\[ T_\omega: \scrD^{2d-k}(M) \to \bbC, \quad \eta \mapsto \int_M \omega \wedge \eta. \]
\end{example}

% \begin{proposition}\label{prop:partial_and_partial_bar_are_differential_operators}
% 	The operators \(\partial\) and \(\overline{\partial}\) satisfy
% 	\[ \partial^2 = 0, \quad \overline{\partial}^2 = 0, \quad \partial \overline{\partial} + \overline{\partial} \partial = 0. \]
% \end{proposition}
% \begin{proof}
% 	\Yang{To be added.}
% \end{proof}
%
% \begin{proposition}\label{prop:holomorphic_pull_back_compatiable_with_A^pq}
% 	Let \(f: M \to N\) be a holomorphic map between complex manifolds.
% 	Then the pull-back of differential forms \(f^*: \calA^{k}_N \to \calA^{k}_M\) satisfies
% 	\[ f^*(\calA^{p,q}_N) \subset \calA^{p,q}_M, \quad f^* \circ \partial_N = \partial_M \circ f^*, \quad f^* \circ \overline{\partial}_N = \overline{\partial}_M \circ f^*. \]
% \end{proposition}
% \begin{proof}
% 	\Yang{To be added.}
% \end{proof}

% \Yang{The following need to checked.}
% \paragraph{Topological vector space of forms with compact support}
% Let \(M\) be a complex manifold of complex dimension \(d\).
% Given a differential form \(\omega \in \calA^{k}(M)\) with compact support, for any compact subset \(K \subset M\) and non-negative integer \(m\), we can define a seminorm
% \[ p_{K,m}(\omega) = \sup_{x \in K} \max_{|\alpha| \leq m} |D^{\alpha} \omega(x)|. \]
% Here, \(\alpha = (\alpha_1, \ldots, \alpha_{2d})\) is a multi-index, and \(D^{\alpha} = \frac{\partial^{|\alpha|}}{\partial x_1^{\alpha_1} \cdots \partial x_{2d}^{\alpha_{2d}}}\) in local real coordinates \((x_1, \ldots, x_{2d})\).
% The collection of these seminorms endows the space of compactly supported smooth \(k\)-forms on \(M\) with a locally convex topology.
%
%
% \begin{definition}\label{def:currents}
% 	A \emph{current} of degree \(k\) on a complex manifold \(M\) is a continuous linear functional on the space of compactly supported smooth \((2d - k)\)-forms on \(M\):
% 	\[ T: \calA^{2d - k}_c(M) \to \bbC. \]
% 	The space of currents of degree \(k\) on \(M\) is denoted by \(\calD_k(M)\).
% 	\Yang{To be revised.}
% \end{definition}



\subsection{Pluripotential theory}



\subsection{Bifurcation measures and currents}


